\documentclass[10pt,a4paper]{amsart}
\usepackage[margin=19mm]{geometry}
\usepackage{amsmath,amssymb,mathtools}
\usepackage[scr=rsfs,cal=euler]{mathalfa}
\usepackage{xfrac}
\usepackage[unicode,hidelinks,bookmarks=false]{hyperref}
\newcommand{\ie}{i.e.\ }
\newcommand{\cf}{cf.\ }
\newcommand{\resp}{respectively\ }
\newcommand{\Subsection}{\subsection}
\providecommand{\nobreakdash}{\nobreak\hbox{-}}
\setlength{\emergencystretch}{2em}
\multlinegap0pt
\usepackage{bm}
\usepackage{bbm}
\usepackage{dsfont}
\usepackage{xspace}
\usepackage{cancel}
\usepackage{mathtools}
\usepackage{amsthm}
\makeatletter
\@ifundefined{theorem}{\newtheorem{theorem}{Theorem}}{}
\@ifundefined{lemma}{\newtheorem{lemma}[theorem]{Lemma}}{}
\makeatother
\title[Unavoidable induced subgraphs in graphs with complete bipart]{Unavoidable induced subgraphs in graphs with complete bipartite induced minors}
\author{Maria Chudnovsky, Meike Hatzel, Tuukka Korhonen, Nicolas Trotignon, Sebastian Wiederrecht}
\date{Source: arXiv:2405.01879v2 (CC BY 4.0)}
\begin{document}
\maketitle

\noindent\emph{Source and licence.} Unavoidable induced subgraphs in graphs with complete bipartite induced minors. Version arXiv:2405.01879v2 (CC BY 4.0), distributed under the Creative Commons Attribution 4.0 International licence, as recorded in the arXiv licence metadata for this exact version and shown on the abstract page https://arxiv.org/abs/2405.01879v2. Full rights notice: licenses/arxiv-2405.01879-CC-BY-4.0.txt.\par\smallskip

\noindent\textit{Editorial scope.} Counted unit: the Lemma whose source label is \texttt{l:OnePath} in the article (source0.tex lines 595--602, source label \texttt{l:OnePath}), delivered together with the article's own complete original proof of it (source0.tex lines 604--649, one \texttt{proof} environment of 46 lines). No mathematics is written, repaired or re-derived: statement and proof are the article's own text line for line. Required context retained uncounted: l:type (lines 544--550). Every definition, display and label of those lines is retained unchanged. Omitted from this package and not redistributed: 1--543, 551--594, 603--603, 650--1353. Nothing inside a retained excerpt is omitted or altered: every retained body is delivered exactly as the article's lines are written, so no omission receipt of any kind accompanies this package. Citations: the retained excerpts carry no citation command at all, so no bibliography entry of the article is transcribed and no reference is redistributed. Reference adaptation: none was needed and none was made; every reference inside the delivered unit resolves inside it, so no unresolved reference or citation is produced. Printed slips kept literal: no printed word, symbol, hypothesis or number of the counted statement or of its proof was altered. Layout and class: the article's own class is \texttt{article} with packages including setspace, amsthm, amssymb, amstext, amsmath, latexsym, graphicx, hyperref, cleveref, thm-restate, tikz, mathtools, enumerate, paralist; the wrapper is the lane's pinned \texttt{amsart} editorial wrapper at 10pt on A4, which restates the theorem-like environments the retained text needs and adds the article's own macro definitions that the retained text needs (none), with extra wrapper packages (bm, bbm, dsfont, xspace, cancel, mathtools). The delivered numbering is the unit's own. Assets: no retained span contains a figure environment, a graphics-inclusion command, a PDF-page inclusion, a table or a TikZ picture; no image file is copied into this package, the package declares no asset dependency and no image file is redistributed with it. Licence: version arXiv:2405.01879v2 of this article is distributed under the Creative Commons Attribution 4.0 International licence, as recorded in the arXiv licence metadata for this exact version and shown on the abstract page https://arxiv.org/abs/2405.01879v2. A full rights notice is delivered at licenses/arxiv-2405.01879-CC-BY-4.0.txt. Characters: every retained excerpt is pure ASCII, so the delivered engine, XeLaTeX with its default Latin Modern text font, reports no missing character.
\begin{lemma}
  \label{l:type}
  Let $G$ be a graph and $A$, $X$, $Y$ and $Z$ be disjoint sets of
  vertices of $G$.  If $A$ is connected and $A$ sees $X$, $Y$ and $Z$,
  then $A$ is of type path, of type claw, or of type triangle with
  respect to $X$, $Y$ and $Z$.
\end{lemma}

\begin{lemma}
  \label{l:OnePath}
  Let $G$ be a 3PC-free graph and $A$, $B$, $X$, $Y$ and $Z$ be
  disjoint connected subsets of $V(G)$. If $X$, $Y$ and $Z$ are
  pairwise anticomplete, $A$ and $B$ are anticomplete to each other,
  and each of $A$ and $B$ sees each of $X$, $Y$ and $Z$, then $A$ and $B$ are
  of type path with respect to $X$, $Y$ and $Z$.
\end{lemma}

\begin{proof}
  Suppose for a contradiction that $A$ (the argument for $B$ is the same by symmetry) is not of type path with
  respect to $X$, $Y$ and $Z$.  Hence, by Lemma~\ref{l:type}, $A$ we
  may consider the two cases below.

  \noindent{\textbf{Case 1:}} $A$ is of type claw (and not of type path).  So, there
  exists a vertex $a\in A$ and three paths $P = a\dots x$,
  $Q = a\dots y$ and $R= a\dots z$ like in the definition of type
  claw.  Moreover, $a\neq x$, $a\neq y$ and $a\neq z$ for otherwise,
  $A$ would be of type path.

  If $B$ is of type claw, then there exist a vertex $b\in B$ and
  three paths $P' = b\dots x'$, $Q' = b\dots y'$ and $R'= b\dots z'$
  like in the definition of type claw. Consider a shortest path $P''$
  from $x$ to $x'$ with interior in $X$, and let $Q'' = y\dots y'$ and
  $R'' = z \dots z'$ be defined similarly through $Y$ and $Z$.  The
  nine paths $P$, $Q$, $R$, $P'$, $Q'$, $R'$, $P''$, $Q''$ and $R''$
  form a theta from $a$ to $b$, a contradiction.

  If $B$ is of type triangle, a similar contradiction is found because
  of the existence of a pyramid in $G$. 

  So $B$ is not of type claw or triangle.  By Lemma~\ref{l:type}, $B$
  is of type path, and up to symmetry we suppose that it is centered
  at $Y$.  Let $P' = x' \dots z'$ be a path like in the definition of
  type path. Consider a shortest path $P''$ from $x$ to $x'$ with
  interior in $X$, and let $R'' = z \dots z'$ be defined similarly
  through $Z$.  Let $Q''=b\dots b'$ be a shortest path in $Y$ such
  that $by\in E(G)$ and $b'$ sees~$P'$.  Let $a'$ be the neighbor of
  $b'$ in $P'$ closest to $x'$ along $P'$. Let $a''$ be the neighbor
  of $b'$ in $P'$ closest to $z'$ along $P'$.  If $a'=a''$, then $B$
  is of type claw, a contradiction.  If $a'a''\in E(G)$, then the
  seven paths $P$, $Q$, $R$, $P'$, $P''$, $Q''$ and $R''$ form a
  pyramid from $a$ to $b'a'a''$, a contradiction.  Hence $a\neq a'$
  and $aa'\notin E(G)$.  So the paths $P$, $Q$, $R$, $x'P'a'$,
  $a''P'z'$, $P''$, $Q''$ and $R''$ form a theta from $a$ to $b'$, a
  contradiction.


  \noindent{\textbf{Case 2:}} $A$ is of type triangle.

  The proof is almost the same as in Case~1, so we just sketch it.  If
  $B$ is of type claw, then $G$ contains a pyramid. If $B$ is of type
  triangle, then $G$ contains a prism. And if $B$ is of type path,
  then $G$ contains a prism or a pyramid.   
\end{proof}

\end{document}
